\chapter*{A field guide to the next problem}
\markboth{A field guide to the next problem}{A field guide to the next problem}
\addcontentsline{toc}{chapter}{A field guide to the next problem}
When a problem looks unfamiliar, first translate it into a question you can ask precisely. A good tool is often a change of language rather than a new formula.

\begin{description}
\item[A finite sum.] Write the first and last terms. Try a difference $F(k+1)-F(k)$, a double count, or the binomial basis. Do not lose an endpoint in a change of index.
\item[A complex expression.] Decide whether addition (coordinates) or multiplication (modulus and argument) dominates. Arguments are defined modulo $2\pi$; roots require all possible arguments.
\item[A polynomial identity.] Bound the degree of the difference, then exhibit enough distinct roots. If the coefficient domain is not a field, revisit the root-counting argument before using it.
\item[An inequality.] Identify the sign of every denominator. Homogeneity may remove a scale. Find the equality case before calculating: it often reveals the correct square or tangent.
\item[A sequence.] Prove it stays in its domain. Search for monotonicity, an invariant interval or a Cauchy estimate. Solving the fixed-point equation identifies possible limits, not convergence.
\item[A series.] Test the term first, then positivity or absolute convergence. A term tending to zero is necessary and insufficient. For a decimal, control the tail; for a rearrangement, justify the operation.
\item[An integral.] Distinguish a finite interval from an improper endpoint. Apply substitutions to truncated intervals first and justify the limiting step. A finite Taylor formula with a bounded remainder often avoids a more difficult interchange theorem.
\item[A matrix.] Ask for the kernel, image and invariant subspaces before expanding determinants. Similarity changes coordinates, not the map. An orthogonal projection is a best approximation; its residual is a certificate.
\item[A probability.] Specify outcomes, probabilities and the information revealed. Count ordered or unordered objects consistently. Add expectations freely; multiply probabilities only under the right independence hypothesis.
\end{description}

\section*{How to read a proof actively}
Cover the next line and predict what would make the current line useful. If you cannot, ask for a small example and a boundary case. At the end, remove one hypothesis and try to break the conclusion. This is not vandalism: it is how you learn which wall holds up the roof.

For a second reading, revisit the eight workshops. Each has a complete worked example, a mission with a staged hint, a solution and a warning about a plausible wrong argument. Try to explain the warning to an imaginary friend who is bright, impatient and very fond of dividing by zero. If your friend survives the explanation, try a real one.

\section*{Computers as experimental partners}
Use exact checks where possible, draw examples, and compare numerical approximations against proved error bounds. The calculations can detect a mistake in an identity or an implementation; they cannot establish a universal theorem from finitely many cases. Floating-point output also carries rounding error. A symbolic equality, a truncation bound and a plotted approximation are three different kinds of evidence; keep their labels attached.
